\pdfvariable suppressoptionalinfo 512
\pdfvariable trailerid {[<C3692026090400000000000000000000><C3692026090400000000000000000000>]}
\documentclass[10pt]{article}
\usepackage[margin=0.72in]{geometry}
\usepackage{amsmath,amssymb,amsthm,booktabs,array,float,microtype,xurl}
\usepackage{unicode-math}
\setmainfont{Latin Modern Roman}
\setmathfont{Latin Modern Math}
\providecommand{\CRevisionRound}{2}
\newtheorem{theorem}{Theorem}
\newtheorem{proposition}{Proposition}
\newtheorem{lemma}{Lemma}
\newcommand{\Fp}{\mathbb F_p}
\newcommand{\Gal}{\operatorname{Gal}}
\newcommand{\Fix}{\operatorname{Fix}}
\newcommand{\Spec}{\operatorname{Spec}}
\newcommand{\Scope}{\texttt{NO\_BAD\_EULER\_OR\_ROOT\_NUMBER}}
\title{The Complete $S_4$ Chebotarev--Frobenius Permutation Atlas\\
for the Root Scheme of $x^4-x-1$}
\author{HCS-C369 / HEN-O353}
\date{4 September 2026}
\begin{document}
\maketitle

\begin{abstract}
The integral quartic $f(x)=x^4-x-1$ defines an intrinsic prime-indexed
family of finite dynamical systems: arithmetic Frobenius
$F_p(\alpha)=\alpha^p$ permutes the four geometric roots at every good
prime.  We prove that $\operatorname{disc}(f)=-283$ and
$\Gal(f/\mathbb Q)=S_4$, identify factor degrees with primitive cycle
lengths, and close the fixed-point and finite-determinant identities on
every good fiber.  All five $S_4$ cycle types are exhibited.  The final
round adds their Chebotarev densities, the non-\'etale boundary at $p=283$,
and the exact fiberwise Koopman classification.  The construction is a
source-arithmetic candidate, not a cross-prime target determinant.  The
universal finite-permutation trace/determinant mechanism remains owned by
the prior workspace package HCS-C12A; our owner is the complete
$x^4-x-1$-specific atlas and its executable convention lock.
\end{abstract}

\section{Quartic arithmetic and the $S_4$ proof}
Put
\begin{equation}\label{eq:f}
 f(x)=x^4-x-1,\qquad
 \mathcal X=\Spec\mathbb Z[x]/(f).
\end{equation}
For a rational prime $p$, write $X_p$ for the geometric roots of the
reduction of $f$ in $\overline{\Fp}$.  We use \emph{arithmetic Frobenius}
\begin{equation}\label{eq:frob}
 F_p:X_p\longrightarrow X_p,\qquad \alpha\longmapsto\alpha^p.
\end{equation}
Geometric Frobenius is $F_p^{-1}$.  The two conventions have identical
cycle lengths, fixed counts, and determinants; this paper never conflates
the maps, only their common cycle partition.

\begin{theorem}[Arithmetic and Galois closure]\label{thm:galois}
The discriminant of $f$ is $-283$.  The polynomial is irreducible over
$\mathbb Q$ and its splitting field $L$ has
\[
 \Gal(L/\mathbb Q)\cong S_4.
\]
Every prime $p\ne283$ is unramified and $\mathcal X_{\Fp}$ is a finite
\'etale scheme of degree four.  If
$\lambda_p=(d_1,\ldots,d_k)$ is the multiset of degrees of the distinct
irreducible factors of $f\bmod p$, then $\lambda_p$ is exactly the cycle
partition of $F_p$ on $X_p$.
\end{theorem}

\begin{proof}
For $x^4+qx+r$ the discriminant is $256r^3-27q^4$; inserting $q=r=-1$
gives $-256-27=-283$.  Modulo $2$,
\[
 f=x^4+x+1.
\]
It has no root, and the only monic irreducible quadratic over $\mathbb F_2$
is $x^2+x+1$, whose square is $x^4+x^2+1$.  Hence $f\bmod2$ is
irreducible.  Gauss's lemma makes $f$ irreducible over $\mathbb Q$, so its
Galois group $G\le S_4$ is transitive.  Dedekind's theorem~\cite{Serre,Neukirch} at the
unramified prime $2$ supplies a $4$-cycle.  At $p=7$,
\begin{equation}\label{eq:p7}
 f\equiv(x-3)(x^3+3x^2+2x-2)\pmod 7,
\end{equation}
and the cubic factor has no root, so $G$ also contains a $3$-cycle.
Consequently $12$ divides $|G|$, while $|G|$ divides $24$.  A subgroup of
order $12$ in $S_4$ is the index-two sign kernel $A_4$, but a $4$-cycle is
odd; equivalently, the discriminant is not a square.  Therefore $G=S_4$.

The discriminant criterion makes precisely $p=283$ exceptional.
At every other prime, the roots are distinct.  An irreducible factor of
degree $d$ has its roots in one orbit of $\alpha\mapsto\alpha^p$ of length
$d$, and every orbit arises this way~\cite{Lidl}.  This proves the dictionary.
\end{proof}

\begin{proposition}[Five exact witnesses]\label{prop:witnesses}
Every partition of four occurs among the good fibers.  In addition to
\eqref{eq:p7}, one has
\begin{align*}
 p=2 &: & f&=x^4+x+1 &&(4),\\
 p=17&: & f&=(x+2)(x+5)(x^2-7x+5) &&(2+1+1),\\
 p=71&: & f&=(x^2+15x-20)(x^2-15x+32) &&(2+2),\\
 p=83&: & f&=(x+3)(x+7)(x+14)(x-24) &&(1+1+1+1),
\end{align*}
where every equality is in the corresponding residue field and the final
parenthesis is the Frobenius cycle partition.
\end{proposition}

\begin{proof}
Multiplication gives $f$ modulo the displayed prime.  The linear factors
are explicit; the remaining quadratic or cubic factors have no residue-field
root.  The mod-$2$ case was proved in Theorem~\ref{thm:galois}.
\end{proof}

\ifnum\CRevisionRound>0
\section{All-iterate orbit and determinant closure}
HCS-C12A already owns the universal zero-dimensional statement that
Frobenius on a reduced finite fiber is a permutation, its iterate counts are
traces, and its finite zeta is the reciprocal permutation determinant.
C369 does not claim workspace ownership of that universal mechanism.  This
section specializes it to the five $S_4$ classes of $x^4-x-1$ and locks the
factor, fixed, primitive, and determinant conventions in one executable
all-good-prime ledger.
For a good $p$, let $P_p$ be the permutation matrix on
$\mathbb C[X_p]$ defined by $P_pe_\alpha=e_{F_p(\alpha)}$, and set
\[
 N_p(r)=\#\Fix(F_p^r).
\]

\begin{theorem}[Fixed, primitive, and determinant identities]\label{thm:zeta}
Let $p\ne283$, let $\lambda_p=(d_1,\ldots,d_k)$, and let $r,n\ge1$.
Then
\begin{equation}\label{eq:fixed}
 N_p(r)=\sum_{j:d_j\mid r}d_j.
\end{equation}
If $E_p(n)$ is the number of points of exact period $n$ and $C_p(n)$ the
number of primitive $n$-cycles, then
\begin{equation}\label{eq:primitive}
 E_p(n)=\sum_{d\mid n}\mu(d)N_p(n/d),
 \qquad C_p(n)=\frac{E_p(n)}n.
\end{equation}
Finally, in $\mathbb Q[[u]]$ and therefore as a rational function,
\begin{equation}\label{eq:zeta}
 \begin{split}
 Z_p(u)&:=\exp\!\left(\sum_{r\ge1}\frac{N_p(r)}r u^r\right)\\
 &=\det(I-uP_p)^{-1}
 =\prod_{j=1}^k(1-u^{d_j})^{-1}.
 \end{split}
\end{equation}
Thus the factor degrees are not merely a fixed-point statistic: they are
the complete primitive-cycle ledger on the fiber.
\end{theorem}

\begin{proof}
A $d$-cycle contributes all its $d$ points to $\Fix(F_p^r)$ exactly when
$d\mid r$, proving \eqref{eq:fixed}.  The relation
$N_p(r)=\sum_{n\mid r}E_p(n)$ and M\"obius inversion give
\eqref{eq:primitive}; division by $n$ groups exact-period points into
cycles.  On a $d$-cycle block $C_d$,
$\det(I-uC_d)=1-u^d$.  Multiplying blocks proves the determinant equality,
while
\[
 -\log(1-u^d)=\sum_{m\ge1}\frac{u^{dm}}m
 =\sum_{m\ge1}\frac d{dm}u^{dm}
\]
and \eqref{eq:fixed} prove the exponential equality.
\end{proof}

\begin{table}[H]
\centering
\small
\begin{tabular}{@{}lrrrrr@{}}
\toprule
type & $1+1+1+1$ & $2+1+1$ & $2+2$ & $3+1$ & $4$\\
\midrule
good primes $\le10^4$ & 43 & 306 & 147 & 411 & 321\\
$S_4$ class size & 1 & 6 & 3 & 8 & 6\\
\bottomrule
\end{tabular}
\caption{Exact finite regression counts.  The 1,228 entries sum across the
row; these counts are not used to prove an asymptotic density.}
\label{tab:counts}
\end{table}

The canonical evidence enumerates all $1{,}229$ primes at most $10^4$,
removes the unique bad prime $283$, and checks \eqref{eq:fixed}--\eqref{eq:zeta}
for $1\le r\le12$ on all $1{,}228$ good fibers: $14{,}736$ prime--iterate
cells.  A separate factorization backend recomputes every partition.  Exact
symbolic and hostile-mutation lanes test the proof anchors and schema.
Finite enumeration is an implementation receipt; Theorems
\ref{thm:galois} and \ref{thm:zeta} establish the quartic specialization
for all good primes and all iterates.
\fi

\ifnum\CRevisionRound>1
\section{Chebotarev density, bad fiber, and operator boundary}
This round completes the quartic-specific density, bad-fiber, and route
boundary.

\begin{theorem}[Complete density and boundary atlas]\label{thm:density}
The natural densities of good primes with the five factor partitions are
\begin{center}
\begin{tabular}{@{}lccccc@{}}
\toprule
$\lambda_p$ & $1+1+1+1$ & $2+1+1$ & $2+2$ & $3+1$ & $4$\\
\midrule
density & $1/24$ & $1/4$ & $1/8$ & $1/3$ & $1/4$\\
\bottomrule
\end{tabular}
\end{center}
At the excluded prime,
\begin{equation}\label{eq:bad}
 f\equiv(x-115)(x-93)^2(x+18)\pmod {283},
 \qquad \gcd(f,f')=x-93.
\end{equation}
Hence $\mathcal X_{\mathbb F_{283}}$ is non-\'etale; no four-distinct-point
permutation atlas is asserted there.
\end{theorem}

\begin{proof}
The conjugacy classes of $S_4$ with cycle types
$1+1+1+1,2+1+1,2+2,3+1,4$ have sizes $1,6,3,8,6$.
By Theorem~\ref{thm:galois}, Dedekind's dictionary identifies these classes
with the factor partitions.  Chebotarev's density theorem divides each
class size by $|S_4|=24$~\cite{Neukirch}.  Multiplication verifies \eqref{eq:bad}; the common
root $93$ of $f$ and $f'$ proves nonreducedness.
\end{proof}

\begin{proposition}[Natural fiber quantization]\label{prop:koopman}
For each good $p$, $P_p$ is a canonical unitary on the four-dimensional
Hilbert space $\ell^2(X_p)$.  It is self-adjoint exactly for types
$1+1+1+1$, $2+1+1$, and $2+2$.  For types $3+1$ and $4$ it remains unitary
but is not self-adjoint.
\end{proposition}

\begin{proof}
$P_p$ permutes an orthonormal basis, so $P_p^*=P_p^{-1}$.  It is
self-adjoint precisely when $P_p=P_p^{-1}$, or $F_p^2=I$; this is equivalent
to every cycle having length at most two.
\end{proof}

\section{Route decision and nonclaims}
The integral scheme supplies prime fibers, prime-power Frobenius iterates,
Chebotarev classes, and source-derived primitive weights.  It therefore
earns $A0_{\rm STRUCTURAL\_ARITHMETIC\_RELATION}$.  The exact rational
finite-fiber determinant in \eqref{eq:zeta} earns the deliberately local
$A1_{\rm PASS\_ANALYTIC}$ as an exact applicability result, not as new
ownership of the universal mechanism.  Proposition~\ref{prop:koopman} gives
$A4_{\rm NATURAL\_QUANTIZATION}$.

The prime $p$ indexes a different fiber; it is not a time step of one
autonomous system.  We construct no cross-prime Fredholm direct sum, no
determinant-class regularization, and no target Euler product.  There is no
target continuation, functional equation, divisor, zero fit, or counting
law.  Thus $A2_{\rm FAIL}$ and $A3_{\rm FAIL}$.  The overall verdict is
\texttt{ROUTE\_A\_ARITHMETIC\_CANDIDATE}, under the literal scope
\Scope.  Fiber permutation unitaries are not a Hilbert--P\'olya operator,
and Route B is not invoked.
\fi

\begin{thebibliography}{9}\small
\bibitem{Serre} J.-P. Serre,
\emph{Topics in Galois Theory}, second edition, A K Peters, 2008.
\bibitem{Neukirch} J. Neukirch,
\emph{Algebraic Number Theory}, Springer, 1999.
\bibitem{Lidl} R. Lidl and H. Niederreiter,
\emph{Finite Fields}, second edition, Cambridge University Press, 1997.
\end{thebibliography}
\end{document}
